\section{A physical minimum in a declared architecture}
\label{sec:pa-physical-minimum}

An upper bound on a relaxed record divergence is not a realizable upper bound
on the least dissipation compatible with a measurement. This section gives
an explicit physical construction and an exact optimization within one
architecture. The architecture assumption is substantial: translation-invariant
Markov dynamics on a ring, with the first Fourier mode observed. It is not
inferred merely from two measured correlations.

Let $n\geq3$ and let $k_j\geq0$ be the rate of the displacement
$x\mapsto x+j\pmod n$. Assume the displacement support is irreducible.
The stationary distribution is uniform. Write the first Fourier eigenvalue as
$-\gamma+i\omega$, where $\gamma>0$, so that
\begin{align}
 \gamma&=\sum_{j=1}^{n-1}k_j\{1-\cos(2\pi j/n)\},\\
 \omega&=\sum_{j=1}^{n-1}k_j\sin(2\pi j/n),\\
 \sigma&=\sum_{j=1}^{\lfloor(n-1)/2\rfloor}
       (k_j-k_{n-j})\log(k_j/k_{n-j}).
 \label{eq:pa-circulant-epr}
\end{align}
Terms with both rates zero contribute zero; a positive rate with zero reverse
rate contributes infinite entropy production. A self-inverse displacement
$j=n/2$ contributes no entropy. Entropy is dimensionless per unit time;
converting it to physical power requires the thermodynamic embedding.

The Fourier formulas, the general delocalized-mode inequality
$\sigma\geq\omega^2/\gamma$, and its saturation by circle drift--diffusion
are established results, discussed explicitly by Gu~\cite{rfGu2026}.
The finite-$n$ minimization below supplies a more restrictive, architecture-specific
answer. No historical-priority claim is made.

\paragraph{Theorem: finite-size circulant minimum.}
Fix $n,\gamma,\omega$. A finite-entropy circulant realization exists if and
only if
\begin{equation}
 \frac{|\omega|}{\gamma}<\cot(\pi/n).
 \label{eq:pa-circulant-feasible}
\end{equation}
Within that class the minimum is
\begin{equation}
 \boxed{\quad
 \sigma_n^{\min}(\gamma,\omega)
 =\frac{2\omega}{\sin(2\pi/n)}
   \operatorname{artanh}\!\left(
       \frac{\omega}{\gamma}\tan(\pi/n)\right).
 \quad}
 \label{eq:pa-circulant-min}
\end{equation}
The expression is even in $\omega$ and is zero at $\omega=0$. The minimizer
uses only the two nearest-neighbor displacements, with rates
\begin{equation}
 k_{\pm}=\frac12\left\{
 \frac{\gamma}{1-\cos(2\pi/n)}
 \mathbin{\pm}\frac{\omega}{\sin(2\pi/n)}\right\}.
 \label{eq:pa-circulant-rates}
\end{equation}
At equality in \eqref{eq:pa-circulant-feasible}, a one-way ring is possible
but has infinite ordinary-reversal entropy production. Above the boundary
there is no such circulant realization, even with infinite entropy.

\paragraph{Proof.}
Pair the displacements with angle $\theta_j=2\pi j/n\in(0,\pi)$.
Set $x_j=\tan(\theta_j/2)$, traffic $s_j=k_j+k_{n-j}$ and
current $d_j=k_j-k_{n-j}$. A nonzero pair contributes
\[
 \gamma_j=s_j(1-\cos\theta_j),\qquad
 \omega_j=d_j\sin\theta_j,\qquad
 \sigma_j=2d_j\operatorname{artanh}(d_j/s_j).
\]
For $u_j=\omega_j/\gamma_j$ this is
\[
 \sigma_j=\gamma_j F_{x_j}(u_j),\qquad
 F_x(u)=\frac{u(1+x^2)}{x}\operatorname{artanh}(xu),
 \qquad |xu|<1.
\]
The function is even in $u$, convex because
$F_x''(u)=2(1+x^2)/(1-x^2u^2)^2$, and nondecreasing in $x$ at fixed $u$.
For the last assertion, $\operatorname{artanh}(xu)/x$ is nondecreasing
for $u\geq0$, as follows from its positive power series; evenness handles
$u<0$. Since $x_j\geq x_*=\tan(\pi/n)$,
$F_{x_j}(u_j)\geq F_{x_*}(u_j)$. An antipodal displacement, when present,
adds damping with $u_j=0$ and zero entropy and can be included as another
weight. Jensen's inequality with weights $\gamma_j/\gamma$ gives
\[
 \sigma\geq\sum_j\gamma_j F_{x_*}(u_j)
 \geq\gamma F_{x_*}(\omega/\gamma),
\]
which is \eqref{eq:pa-circulant-min}. If a reverse rate vanishes the same
inequality holds with infinite left side. Also,
$|\omega_j|/\gamma_j\leq1/x_j\leq1/x_*$, strictly below the last
boundary for a finite-entropy nonzero oscillation. This proves necessity
of \eqref{eq:pa-circulant-feasible}. The positive rates
\eqref{eq:pa-circulant-rates} prove sufficiency and attain equality.
\hfill$\square$

For fixed nonzero $\omega$, the minimum decreases strictly as feasible $n$
increases, and
\begin{equation}
 \sigma_n^{\min}>\frac{\omega^2}{\gamma},\qquad
 \lim_{n\to\infty}\sigma_n^{\min}=\frac{\omega^2}{\gamma}.
 \label{eq:pa-circulant-limit}
\end{equation}
This follows from the same monotonicity of $F_x$ and $x=\tan(\pi/n)\downarrow0$.
Equivalently,
\[
 \frac{\sigma_n^{\min}}{\omega^2/\gamma}
 =\sec^2(\pi/n)\,
   \frac{\operatorname{artanh}(r_n)}{r_n},\qquad
 r_n=(\omega/\gamma)\tan(\pi/n).
\]
The finite-size overhead is physical, not a sampling-confidence penalty.

\paragraph{A complete correlation law with explicit realizations.}
For $0<A\leq1/2$, the bounded detector
$f(x)=\sqrt{2A}\cos(2\pi x/n)$ has
\begin{equation}
 C(t)=\mathbb E[f(X_0)f(X_t)]
     =A e^{-\gamma t}\cos(\omega t).
 \label{eq:pa-mode-correlation}
\end{equation}
Every ring constructed above realizes this same law for all $t\geq0$.
The circle diffusion
\[
 d\phi_t=\omega\,dt+\sqrt{2\gamma}\,dW_t\pmod{2\pi},\qquad
 f(\phi)=\sqrt{2A}\cos\phi
\]
has uniform stationary distribution, exactly the same correlation, and
$\sigma=\omega^2/\gamma$. This is an attainable physical value, unlike
an upper bound on a record-divergence relaxation. It provides an upper
bound on the \emph{infimum} of entropy among all compatible models in
any larger class containing this diffusion. It is not an upper bound on
the entropy of the unknown system. In the union of circulant finite-state
architectures, it is an infimum approached by finite rings.

For the earlier calibrated example $C(0)=1/2$, $C(1)=0.38$ and
$C(2)=0.18$, impose the single-mode model and write
$z=C(1)/A$, $w=C(2)/A$. Then
\[
 e^{-2\gamma}=2z^2-w=497/625,
 \quad \gamma=0.1145808118\ldots,
 \quad |\omega|_{\rm principal}=0.5504934894\ldots.
\]
The smallest feasible ring has $n=16$. Selected exact-form physical
minima, evaluated numerically, are
\begin{center}
\begin{tabular}{r|rrrrr|r}
$n$&16&20&40&100&1000&circle\\\hline
$\sigma_n^{\min}$&5.44695&3.55736&2.80013&2.66781&2.64502&2.64480
\end{tabular}
\end{center}
These numbers sharpen the declared ring problem, not the unrestricted
inverse problem. A hidden model outside this architecture can have the
same measured moments. Eigenvalue-only entropy bounds require particular
care outside the delocalized class; known counterexamples preclude replacing
this declaration with a universal claim~\cite{rfShiraishi2026}.

\paragraph{An explicit failure outside the architecture.}
Even the complete scalar correlation function does not by itself justify the
circulant minimum. Take a four-state nearest-neighbor ring with damping
$\gamma-\kappa$ and frequency $\omega$, where
$0<\kappa<\gamma-|\omega|$. Add a fifth state, called the hub.
From every ring state jump to the hub at rate $\kappa$; from the hub jump
to each ring state at rate $\beta/4$, where
$\beta=A\kappa/(1-A)$ and $0<A<1$. The stationary masses are
$\pi_i=A/4$ on the ring and $\pi_\mathrm{hub}=1-A$.
The hub edges obey detailed balance. The detector is
$f=(1,1,-1,-1,0)$, so $\mathbb E f=0$ and $\mathbb E f^2=A$.
The ring Fourier subspace has zero hub component and its damping increases
by $\kappa$, giving exactly
\[
 C(t)=A e^{-\gamma t}\cos(\omega t),\qquad
 \sigma=2A\omega\operatorname{artanh}\!\left(
                       \frac{\omega}{\gamma-\kappa}\right).
\]
The chain is irreducible and every nonzero edge has a positive reverse.
At $A=0.1$, $\gamma=1$, $\omega=0.5$, $\kappa=0.01$, the entropy is
$0.0556063004\ldots$, below the circle value $0.25$, despite identical
calibration and identical scalar covariance at every lag. Direct generator
calculations reproduce both values. Thus a circle construction provides a
compatible physical upper bound on the infimum; it does not prove the
unrestricted infimum equals that construction. The detector law beyond
second-order correlations may distinguish these architectures.

\section{Sampling can hide arbitrarily large physical dissipation}
\label{sec:pa-entropy-aliasing}

The same realizable example exposes a second limitation. Even full spatial
resolution at discrete times need not bound physical dissipation. This is a
thermodynamic consequence of ordinary sampling aliasing, closely related to
continuous-time system aliasing~\cite{paYue2021}.

\paragraph{Exact grid aliasing.}
Fix $\gamma>0$ and a sampling interval $\Delta>0$. The circle transition
kernel at time $\Delta$ depends on $\omega$ only through
$\omega\Delta\pmod{2\pi}$. Therefore every finite or infinite
stationary record sampled on that grid has exactly the same probability
law under
\begin{equation}
 \omega_k=\omega+2\pi k/\Delta,\qquad k\in\mathbb Z.
 \label{eq:pa-frequency-alias}
\end{equation}
This statement holds even if the full circle state is read at each instant.
Any detector or subsequent processing retains the equivalence. Yet
$\sigma_k=\omega_k^2/\gamma\to\infty$. Thus the known two-parameter
architecture and arbitrarily many grid records do not supply a finite
upper bound on dissipation without an additional frequency restriction or
access to intervening motion.

For the scalar covariance, $C(0),C(\Delta),C(2\Delta)$ identify
$\gamma$ and the principal value
\[
 \omega_0=\Delta^{-1}\arccos\!\left[
 \frac{C(\Delta)/A}{\sqrt{2(C(\Delta)/A)^2-C(2\Delta)/A}}\right]
\]
when the specified single-mode model is feasible. The frequencies
$\pm\omega_0+2\pi k/\Delta$ remain possible. Their smallest
circle entropy is $\omega_0^2/\gamma$; the supremum is infinite.
These conclusions concern the single-mode class and exact moments.

\paragraph{Exact identification need not give a stable entropy upper bound.}
If $\gamma$ and $A$ are known, three positive lags with pairwise irrational
ratios determine $|\omega|$ from exact cosine correlations. Indeed,
$\cos(\omega' t_j)=\cos(\omega t_j)$ implies
$\omega't_j=s_j\omega t_j+2\pi k_j$, with $s_j\in\{-1,1\}$.
Among three signs, two agree. An irrational ratio between the associated
lags forces both integers to vanish, and hence $\omega'=\omega$ or
$\omega'=-\omega$. Two such lags do not always suffice: at
$t_1=1,t_2=\sqrt2$, the distinct positive frequencies
$\pi(\sqrt2-1)$ and $\pi(\sqrt2+1)$ give identical cosines.

Exact uniqueness does not survive as a finite-precision upper bound.
Fix a finite menu of positive observation gaps, a fixed finite number of
intervals, and $\gamma>0$. For any reference $\omega$, there is a sequence
$|\omega_k|\to\infty$ whose entire sampled-state record distributions
converge in total variation to the reference distribution. This includes
adaptive choices from the finite gap menu with a bounded number of steps.

To prove the statement, write the gaps as $t_1,\ldots,t_m$ and use
simultaneous Diophantine approximation to choose arbitrarily large integers
$q$ such that the distances from $q t_j/t_1$ to integers approach zero.
If an exact common period exists, its multiples give the conclusion
immediately; otherwise the approximating denominators are unbounded.
Set $\omega_k-\omega=2\pi q/t_1$. All the wrapped phase shifts then
approach zero. For a gap $t$, let $\delta_t$ be its phase shift reduced to
$[-\pi,\pi]$. Comparing unwrapped normal increments and then wrapping gives
\begin{equation}
 \operatorname{TV}(K_\omega^t,K_{\omega_k}^t)
 \leq 2\Phi\!\left(\frac{|\delta_t|}{2\sqrt{2\gamma t}}\right)-1
 \leq\frac{|\delta_t|}{2\sqrt{\pi\gamma t}},
 \label{eq:pa-phase-tv}
\end{equation}
uniformly in the starting state. Sequential coupling bounds complete-record
TV by the sum of these transition bounds, or by the number of steps times
the largest menu bound. Independent replication multiplies the sum bound
by the number of records. Stopping and further detector processing cannot
increase it. Meanwhile $\omega_k^2/\gamma$ diverges.

Consequently every positive-TV neighborhood of the accessible law contains
models with unbounded entropy. More statistically, any upper procedure $U$
with uniform coverage $\Pr_\omega\{U\geq\omega^2/\gamma\}\geq1-\alpha$
for this experiment must satisfy
\begin{equation}
 \Pr_\omega\{U=+\infty\}\geq1-\alpha
 \quad\hbox{at every reference }\omega.
 \label{eq:pa-alias-upper-no-go}
\end{equation}
For any finite $M$, choose near-aliases with entropy exceeding $M$ and
record TV tending to zero. Uniform coverage then implies
$\Pr_\omega\{U\geq M\}\geq1-\alpha$; take $M\to\infty$.
This statement assumes a fixed finite gap menu and bounded experiment.
It does not address a continuously distributed random timing design whose
exact realized times are retained, an unbounded observation horizon, or a
separately calibrated frequency cap. Those changes alter access.

The implementation illustrates the near-alias mechanism constructively.
At gaps $1$ and $\sqrt2$, Pell solutions $p^2-2q^2=1$ give shifts
$2\pi q$, with second-gap phase error
$2\pi/(p+\sqrt2q)$. For $\gamma=1$, reference $\omega=1/2$, and
$(p,q)=(665857,470832)$, the alternative entropy is
$8.751688\times10^{12}$ while the two-interval full-state record TV
is at most $1.119197\times10^{-6}$. This is a mathematical model
calculation, not a physically achievable drive-rate claim. A calibrated
bandwidth excludes such alternatives; absent that calibration the model
allows them.

\paragraph{Reproduction and limits.}
\texttt{physical\_minimum.py} reconstructs generators, compares matrix
exponentials with \eqref{eq:pa-mode-correlation}, checks the inequality on
1,000 seeded nonlocal circulant generators, and saves the physical and
aliasing examples. Fourteen producer tests cover realization, nonlocal
jumps, equilibrium, one-way limits, invalid inputs, continuum convergence,
and exact and near aliasing. Floating-point evaluation is not advertised
as an interval certificate; the proofs establish the mathematical claims.
No new measured-data entropy estimate, universal physical minimum, or
historical priority follows from these calculations.
